\documentclass[rep, style=ml]{lexnote}

\title{Sample: Generative Models}
\author{Jane Doe}
\date{\today}

\begin{document}
\maketitle
\tableofcontents

\chapter{Diffusion Models}

\section{Forward Process}

\begin{defn}[Forward Diffusion]
  The forward process gradually adds Gaussian noise to data $\mathbf{x}_0$
  over $T$ timesteps:
  \[
    q(\mathbf{x}_t \mid \mathbf{x}_{t-1})
    = \mathcal{N}\!\bigl(\mathbf{x}_t;\,
      \sqrt{1-\beta_t}\,\mathbf{x}_{t-1},\;
      \beta_t \mathbf{I}\bigr)
  \]
  where $\{\beta_t\}_{t=1}^T$ is the variance schedule.
  （前向过程：逐步向数据添加高斯噪声）
\end{defn}

Using the reparameterization trick with $\alpha_t = 1 - \beta_t$ and
$\bar{\alpha}_t = \prod_{s=1}^{t} \alpha_s$, we can sample $\mathbf{x}_t$
directly from $\mathbf{x}_0$:

\begin{thm}[Closed-Form Sampling]
  \[
    q(\mathbf{x}_t \mid \mathbf{x}_0)
    = \mathcal{N}\!\bigl(\mathbf{x}_t;\,
      \sqrt{\bar{\alpha}_t}\,\mathbf{x}_0,\;
      (1 - \bar{\alpha}_t)\mathbf{I}\bigr)
  \]
  （可直接从 $\mathbf{x}_0$ 采样任意时刻 $\mathbf{x}_t$，无需逐步迭代）
\end{thm}

\section{Reverse Process}

\begin{lem}[Posterior Distribution]
  When conditioned on $\mathbf{x}_0$, the reverse step has a tractable form:
  \[
    q(\mathbf{x}_{t-1} \mid \mathbf{x}_t, \mathbf{x}_0)
    = \mathcal{N}\!\bigl(\mathbf{x}_{t-1};\,
      \tilde{\boldsymbol{\mu}}_t,\;
      \tilde{\beta}_t \mathbf{I}\bigr)
  \]
  where $\tilde{\boldsymbol{\mu}}_t$ depends on both $\mathbf{x}_t$
  and $\mathbf{x}_0$.
\end{lem}

\subsection{Training Objective}

The \keyword{simplified loss} used in DDPM reduces to a denoising objective:

\begin{notebox}[Key Insight]
  Instead of predicting $\mathbf{x}_0$ directly, the model learns to predict
  the noise $\boldsymbol{\epsilon}$ that was added:
  \[
    L_{\text{simple}}
    = \mathbb{E}_{t,\,\mathbf{x}_0,\,\boldsymbol{\epsilon}}\!\Bigl[
      \bigl\|\boldsymbol{\epsilon}
        - \boldsymbol{\epsilon}_\theta(\mathbf{x}_t, t)\bigr\|^2
    \Bigr]
  \]
  （模型预测噪声而非原始数据，训练更稳定）
\end{notebox}

\begin{rem}[Connection to Score Matching]
  The noise prediction network $\boldsymbol{\epsilon}_\theta$ is equivalent
  to learning the score function
  $\nabla_{\mathbf{x}} \log q(\mathbf{x}_t)$, up to a scaling factor.
  This connects DDPM to score-based generative models.
  （噪声预测等价于学习得分函数，与 score-based 模型统一）
\end{rem}

\chapter{Normalizing Flows}

\section{Change of Variables}

\begin{defn}[Normalizing Flow]
  A normalizing flow transforms a simple base distribution $p_z(\mathbf{z})$
  through a sequence of invertible mappings $f_1, f_2, \ldots, f_K$ to
  produce a complex target distribution:
  \[
    \mathbf{x} = f_K \circ f_{K-1} \circ \cdots \circ f_1(\mathbf{z})
  \]
  （正则化流：通过可逆变换将简单分布映射为复杂分布）
\end{defn}

\begin{thm}[Change of Variables Formula]
  If $f$ is an invertible differentiable function, the density of
  $\mathbf{x} = f(\mathbf{z})$ is:
  \[
    p_x(\mathbf{x})
    = p_z\!\bigl(f^{-1}(\mathbf{x})\bigr)\,
      \left|\det \frac{\partial f^{-1}}{\partial \mathbf{x}}\right|
  \]
  （换元公式：变换后的密度由 Jacobian 行列式调整）
\end{thm}

\section{Architectures}

\begin{defbox}[Coupling Layer]
  Split the input $\mathbf{x}$ into two halves $(\mathbf{x}_a, \mathbf{x}_b)$.
  Transform one conditioned on the other:
  $\mathbf{y}_a = \mathbf{x}_a$,\quad
  $\mathbf{y}_b = \mathbf{x}_b \odot \exp(s(\mathbf{x}_a)) + t(\mathbf{x}_a)$.
  The Jacobian is triangular, so its determinant is cheap to compute.
\end{defbox}

\begin{rem}[Practical Consideration]
  Coupling layers are efficient ($O(D)$ Jacobian) but limited in
  expressiveness per layer. Stacking many layers with alternating
  partitions is essential for modeling complex distributions.
\end{rem}

\todo{Add chapter on VAEs}

\end{document}
